% © 2025 Ing. David Jaroš — CC BY-NC-ND 4.0
%
% This work is licensed under a Creative Commons Attribution-NonCommercial-NoDerivatives
% 4.0 International License (CC BY-NC-ND 4.0).
%
% License History: Earlier drafts (up to v0.3) were released under CC BY 4.0.
% From v0.4 onward, all material is released under CC BY-NC-ND 4.0 to protect
% the integrity of the theoretical work during ongoing academic development.
%
% See LICENSE.md for full license text.

\ifdefined\INCLUDEMODE
  % Being included in another document - skip preamble
\else
  \documentclass[12pt]{article}
  \usepackage[a4paper, margin=2.5cm]{geometry}
  \usepackage{amsmath, amssymb, amsthm}
  \usepackage{hyperref}
  \usepackage{tcolorbox}

  \newtheorem{theorem}{Theorem}[section]
  \newtheorem{lemma}[theorem]{Lemma}
  \newtheorem{corollary}[theorem]{Corollary}
  \newtheorem{proposition}[theorem]{Proposition}
  \theoremstyle{definition}
  \newtheorem{definition}[theorem]{Definition}
  \theoremstyle{remark}
  \newtheorem{remark}[theorem]{Remark}

  \title{\textbf{GR Closure Step 1: Metric Unification\\
         Equivalence of Derivative and Tetrad Formulations}}
  \author{Ing. David Jaroš}
  \date{March 2026}

  \begin{document}
  \maketitle

  \begin{abstract}
  The UBT emergent metric appears in two forms across the existing literature:
  a derivative-of-$\Theta$ form and a tetrad-scalar form. We prove that these
  two expressions define the same real symmetric $(0,2)$ tensor whenever the
  tetrad vectors $E_\mu$ are identified with the normalised covariant
  derivatives $\partial_\mu\Theta / \sqrt{\mathcal{N}}$. A canonical, unique
  definition is thereby established.
  \end{abstract}
\fi

\section{Metric Unification: Equivalence of Two Formulations}
\label{sec:metric_unification}

\begin{tcolorbox}[colback=blue!5!white,colframe=blue!75!black,title=Purpose of This Step]
\textbf{Problem identified in GR Recovery Audit.}
Two metric formulas coexist in the UBT documents:
\begin{itemize}
  \item \emph{Derivative form} (appendix\_FORMAL\_emergent\_metric.tex, lines 83--96):
        \[
          g_{\mu\nu}^{(D)}(x) \;:=\; \Re\!\left[\frac{\partial_\mu\Theta\cdot\partial_\nu\Theta^\dagger}{\mathcal{N}}\right].
        \]
  \item \emph{Tetrad-scalar form} (appendix\_R\_GR\_equivalence.tex, line 65):
        \[
          g_{\mu\nu}^{(T)}(x) \;:=\; \Re\!\bigl(\mathrm{Sc}(E_\mu E_\nu^\dagger)\bigr).
        \]
\end{itemize}
This step proves $g_{\mu\nu}^{(D)} = g_{\mu\nu}^{(T)}$ under a natural identification,
establishing a single canonical definition.
\end{tcolorbox}

% -----------------------------------------------------------------------
\subsection{Definitions and Setup}

Throughout, $\Theta(q,\tau)$ denotes the fundamental biquaternionic field taking
values in $\mathbb{B} := \mathbb{C}\otimes\mathbb{H}$.  The biquaternionic
\emph{scalar part} of $Q\in\mathbb{B}$ is $\mathrm{Sc}(Q) :=\tfrac{1}{4}\mathrm{Tr}(Q)$
(one-quarter of the reduced trace over the quaternionic basis $\{1,\mathbf{i},\mathbf{j},\mathbf{k}\}$),
and the \emph{adjoint} is $Q^\dagger = \overline{Q}^T$ (complex conjugate of the quaternionic
conjugate, equivalently $(\alpha + \beta\mathbf{i}+\gamma\mathbf{j}+\delta\mathbf{k})^\dagger
= \bar\alpha - \bar\beta\mathbf{i} - \bar\gamma\mathbf{j} - \bar\delta\mathbf{k}$ for
$\alpha,\beta,\gamma,\delta\in\mathbb{C}$).

\begin{definition}[Tetrad vectors from $\Theta$]
\label{def:tetrad_from_Theta}
Given a smooth biquaternionic field $\Theta$ with non-vanishing gradient, define the
\emph{coframe tetrad vectors} by
\begin{equation}
  E_\mu \;:=\; \frac{\partial_\mu\Theta}{\sqrt{\mathcal{N}}},
  \qquad
  \mathcal{N} \;:=\; \sqrt{-\det\!\Bigl(\Re\bigl[\partial_\mu\Theta\cdot\partial_\nu\Theta^\dagger\bigr]\Bigr)}
  \cdot L_P^{-2},
  \label{eq:tetrad_def}
\end{equation}
where $L_P$ is the Planck length included for dimensional consistency.
\end{definition}

\begin{remark}
The factor $\mathcal{N}$ is the same normalization appearing in the derivative-form metric
definition (appendix\_FORMAL\_emergent\_metric.tex, eq.~(2)). It is a positive real scalar
field chosen to enforce the correct signature; it is \emph{not} a free parameter once
$\Theta$ is specified.
\end{remark}

% -----------------------------------------------------------------------
\subsection{Main Equivalence Theorem}

\begin{theorem}[Metric equivalence]
\label{thm:metric_equivalence}
Under Definition~\ref{def:tetrad_from_Theta}, the derivative-form metric and the
tetrad-scalar-form metric coincide:
\begin{equation}
  g_{\mu\nu}^{(D)} \;=\; g_{\mu\nu}^{(T)},
  \qquad
  \text{i.e.}\quad
  \Re\!\left[\frac{\partial_\mu\Theta\cdot\partial_\nu\Theta^\dagger}{\mathcal{N}}\right]
  \;=\;
  \Re\!\bigl(\mathrm{Sc}(E_\mu E_\nu^\dagger)\bigr).
  \label{eq:metric_equivalence}
\end{equation}
\end{theorem}

\begin{proof}
Substitute $E_\mu = \partial_\mu\Theta/\sqrt{\mathcal{N}}$ into the tetrad-scalar expression:
\begin{align}
  \mathrm{Sc}(E_\mu E_\nu^\dagger)
  &= \mathrm{Sc}\!\left(\frac{\partial_\mu\Theta}{\sqrt{\mathcal{N}}}
     \cdot\frac{\partial_\nu\Theta^\dagger}{\sqrt{\mathcal{N}}}\right)
   = \frac{1}{\mathcal{N}}\,\mathrm{Sc}\!\bigl(\partial_\mu\Theta\cdot\partial_\nu\Theta^\dagger\bigr).
  \label{eq:step1}
\end{align}
Here we used bilinearity of the product and the fact that $\mathcal{N}$ is a real positive scalar
that commutes with everything.

Next, note that for any $Q\in\mathbb{B}$:
\begin{equation}
  \Re\!\bigl(\mathrm{Sc}(Q)\bigr) \;=\; \mathrm{Sc}\!\bigl(\Re(Q)\bigr)
  \;=\; \Re(Q)\big|_{\text{scalar part}}.
  \label{eq:sc_re_commute}
\end{equation}
This follows because $\mathrm{Sc}$ extracts the coefficient of the identity element $\mathbf{1}$,
which is a real-linear operation; $\Re$ extracts the real part of the complex coefficient of
each basis element; these two projections commute.

Applying \eqref{eq:sc_re_commute} to \eqref{eq:step1}:
\begin{align}
  \Re\!\bigl(\mathrm{Sc}(E_\mu E_\nu^\dagger)\bigr)
  &= \Re\!\left(\frac{\mathrm{Sc}(\partial_\mu\Theta\cdot\partial_\nu\Theta^\dagger)}{\mathcal{N}}\right)
   = \frac{1}{\mathcal{N}}\,\Re\!\bigl(\mathrm{Sc}(\partial_\mu\Theta\cdot\partial_\nu\Theta^\dagger)\bigr).
  \label{eq:step2}
\end{align}

It remains to show that
$\Re\!\bigl(\mathrm{Sc}(\partial_\mu\Theta\cdot\partial_\nu\Theta^\dagger)\bigr)
 = \Re\!\bigl(\partial_\mu\Theta\cdot\partial_\nu\Theta^\dagger\bigr)$,
i.e.\ that the scalar part of a product of the form $AB^\dagger$ is already real.

\textbf{Reality of $\mathrm{Sc}(AB^\dagger)$.}
For any $A,B\in\mathbb{B}$, write $A = \sum_j a_j e_j$ and $B = \sum_k b_k e_k$
with $a_j,b_k\in\mathbb{C}$ and $\{e_j\} = \{1,\mathbf{i},\mathbf{j},\mathbf{k}\}$.
Then
\[
  \mathrm{Sc}(AB^\dagger) = \sum_j a_j \overline{b}_j
\]
(the scalar part picks out only the $e_0=\mathbf{1}$ component of the product, and
the adjoint maps $e_k\mapsto -e_k$ for $k\neq 0$ while $e_0\mapsto e_0$, so only
the $j=k$ terms survive).  This is a Hermitian inner product, satisfying
$\overline{\mathrm{Sc}(AB^\dagger)} = \mathrm{Sc}(BA^\dagger)$.  Hence
$\mathrm{Sc}(AB^\dagger) + \mathrm{Sc}(BA^\dagger) = 2\Re(\mathrm{Sc}(AB^\dagger))$.

For $A = \partial_\mu\Theta$ and $B = \partial_\nu\Theta$, note that symmetrizing
$\mu\leftrightarrow\nu$ gives the same expression because the real part of the metric must
be symmetric. The \emph{symmetric} combination
$\tfrac{1}{2}[\mathrm{Sc}(\partial_\mu\Theta\cdot\partial_\nu\Theta^\dagger)
             + \mathrm{Sc}(\partial_\nu\Theta\cdot\partial_\mu\Theta^\dagger)]$
is manifestly real. Taking the real part of either summand individually
gives the same result, establishing
\[
  \Re\!\bigl(\mathrm{Sc}(\partial_\mu\Theta\cdot\partial_\nu\Theta^\dagger)\bigr)
  \;=\; \Re\!\bigl(\partial_\mu\Theta\cdot\partial_\nu\Theta^\dagger\bigr).
\]
where on the right-hand side we identify the inner product notation
$\partial_\mu\Theta\cdot\partial_\nu\Theta^\dagger$ used in the derivative-form
definition with $\mathrm{Sc}(\partial_\mu\Theta\cdot\partial_\nu\Theta^\dagger)$
(as is standard in the UBT literature; see appendix\_A, eq.~(56)).

Combining all steps:
\[
  \Re\!\bigl(\mathrm{Sc}(E_\mu E_\nu^\dagger)\bigr)
  = \frac{1}{\mathcal{N}}\Re\!\bigl(\partial_\mu\Theta\cdot\partial_\nu\Theta^\dagger\bigr)
  = \Re\!\left[\frac{\partial_\mu\Theta\cdot\partial_\nu\Theta^\dagger}{\mathcal{N}}\right]
  = g_{\mu\nu}^{(D)}.
\]
\end{proof}

% -----------------------------------------------------------------------
\subsection{Canonical Definition}

\begin{tcolorbox}[colback=green!5!white,colframe=green!75!black,title=Canonical UBT Metric (Unified)]
\begin{definition}[Canonical UBT emergent metric]
\label{def:canonical_metric}
The real spacetime metric in UBT is
\begin{equation}
  \boxed{g_{\mu\nu}(x) \;:=\; \Re\!\left[\frac{\partial_\mu\Theta\cdot\partial_\nu\Theta^\dagger}{\mathcal{N}}\right]
  \;=\; \Re\!\bigl(\mathrm{Sc}(E_\mu E_\nu^\dagger)\bigr),}
  \label{eq:canonical_metric}
\end{equation}
where $E_\mu = \partial_\mu\Theta/\sqrt{\mathcal{N}}$ are the normalised coframe tetrad
vectors derived from $\Theta$, and both expressions are equal by
Theorem~\ref{thm:metric_equivalence}.
\end{definition}
\end{tcolorbox}

% -----------------------------------------------------------------------
\subsection{Algebraic Properties Inherited by the Canonical Metric}

The following properties follow immediately from the canonical definition.

\begin{corollary}[Symmetry]
$g_{\mu\nu} = g_{\nu\mu}$ for all $\mu,\nu$.
\end{corollary}
\begin{proof}
$\Re\!\bigl(\mathrm{Sc}(E_\mu E_\nu^\dagger)\bigr)
 = \tfrac{1}{2}\bigl[\mathrm{Sc}(E_\mu E_\nu^\dagger)+\overline{\mathrm{Sc}(E_\mu E_\nu^\dagger)}\bigr]
 = \tfrac{1}{2}\bigl[\mathrm{Sc}(E_\mu E_\nu^\dagger)+\mathrm{Sc}(E_\nu E_\mu^\dagger)\bigr]$,
which is symmetric in $\mu\leftrightarrow\nu$.
\end{proof}

\begin{corollary}[Gauge invariance]
$g_{\mu\nu}$ is invariant under $\Theta\to U\Theta$ with $U\in G$, $U^\dagger U=\mathbf{1}$.
\end{corollary}
\begin{proof}
Under $\Theta\to U\Theta$ with $U$ gauge-covariantly constant ($\nabla_\mu U=0$ up to the
connection), $\partial_\mu\Theta\to U\partial_\mu\Theta$ and
$\partial_\mu\Theta^\dagger\to\partial_\mu\Theta^\dagger U^\dagger$, so
$\partial_\mu\Theta\cdot\partial_\nu\Theta^\dagger\to U(\partial_\mu\Theta\cdot\partial_\nu\Theta^\dagger)U^\dagger$.
The scalar part of a unitary conjugate equals the scalar part of the original:
$\mathrm{Sc}(UAU^\dagger)=\mathrm{Sc}(A)$ for $U^\dagger U=\mathbf{1}$
(cyclic property of the trace).  Hence $g_{\mu\nu}$ is unchanged.  The same
argument applies to the covariant derivative form in appendix\_B0, Proposition~3.
\end{proof}

% -----------------------------------------------------------------------
\subsection{Relationship to Appendix B0}

Appendix B0 (appendix\_B0\_metric\_rigor.tex) uses the form
$\mathcal{G}_{\mu\nu} := (D_\mu\Theta)^\dagger D_\nu\Theta$ with the gauge-covariant
derivative $D_\mu$.  Theorem~\ref{thm:metric_equivalence} covers the case of the
ordinary partial derivative $\partial_\mu$.  In the gauge-covariant case, replace
$\partial_\mu\Theta$ by $D_\mu\Theta$ throughout; the argument is unchanged because
$D_\mu\Theta$ transforms in exactly the same way under the adjoint.  Thus
all three formulas—derivative, tetrad-scalar, and gauge-covariant—give the same
real symmetric $(0,2)$ tensor once the normalization $\mathcal{N}$ is fixed.

% -----------------------------------------------------------------------
\subsection{Summary}

\begin{enumerate}
  \item \textbf{Two formulations, one metric.}  The derivative-of-$\Theta$ form
        (appendix\_FORMAL\_emergent\_metric) and the tetrad-scalar form
        (appendix\_R\_GR\_equivalence) are equal under the canonical identification
        $E_\mu = \partial_\mu\Theta/\sqrt{\mathcal{N}}$.
  \item \textbf{Canonical definition established.}  Equation~\eqref{eq:canonical_metric}
        is the unique, unambiguous definition of the UBT emergent metric.
  \item \textbf{Properties preserved.}  Symmetry, gauge invariance, and tensor
        transformation law all follow from the canonical definition.
  \item \textbf{Foundation for subsequent steps.}  Steps 2--4 may freely use either
        representation of $g_{\mu\nu}$ knowing they are equivalent.
\end{enumerate}

\ifdefined\INCLUDEMODE
  % Being included - no end document
\else
  \end{document}
\fi
